% ─────────────────────────────────────────────────────────────────────────────
% Sec. 6  Conclusion.
% ─────────────────────────────────────────────────────────────────────────────
\section{Conclusion}
\label{sec:conclusion}

We presented \method, a training-free method that bounds the memory of high-resolution VLMs by
deciding which tiles to encode before the vision encoder, since token pruning removes tokens only after
they have been computed. On InternVL3.5, \method keeps 95.0\% of the uncompressed accuracy with
5.0$\times$ less memory and is the most accurate method under every memory limit. Its peak memory stays
constant from 7 to 188 tiles, so one GPU serves 41$\times$ more requests than the uncompressed model,
and its behavior carries over from 4B to 14B parameters without re-tuning.

\paragraph{Limitations.}
\method requires a VLM whose vision encoder processes each tile independently, so that the encoder can be
run on a subset of the tiles. All our experiments use InternVL3.5, but the same holds for other tile-based
VLMs: LLaVA-OneVision~\cite{llavaov}, for example, encodes each AnyRes crop separately and adds a
downscaled base image, so \method can be applied to it in the same way, although we have not evaluated it.
Native-resolution encoders~\cite{qwen2vl} attend across the whole image, so encoding only the selected
regions would change the input they were trained on; \method does not cover them.

Improving the recall of the prior for very small objects, evaluating other tile-based VLMs, and
extending the selection to native-resolution encoders and to video are directions for future work.

\section*{LLM Usage Statement}
We used Claude Code (Anthropic) to help write and run experiment code, analyze results, generate tables
and figures, and draft and edit the text of this paper, under the direction of the authors. The authors
reviewed these outputs and take responsibility for all content.
